% !TEX root = ../main.tex
\lecture{4}{2026-10-05}{Folgen}

\section{Folgen und Reihen}
\begin{definition}
    [Folge] (Das n-Glied der Folge) ${\left\{a_n\right\}}_{n \in \mathbb{N}}$ ist eine Funktion a: $\mathbb{N} \rightarrow \mathbb{R}$. Eine Funktionsvorschrift, die Natürlichen Zahlen an Reelen Zahlen zuweist.
\end{definition}

\begin{definition}
    [Reihen] gehören zu einer entsprechenden Folge, wobei sie der n-Summe der Folgenglieder gleichen.
    \begin{equation}
        {\left\{\sum_{i=1}^{n} a_i\right\}}_{n \in \mathbb{N}}
    \end{equation}
\end{definition}

\subsection{Eigenschaften}

\begin{definition}
    Eine Folge ${\left\{a_n\right\}}_{n \in \mathbb{N}}$  heisst:
    \begin{itemize}
        \item \textbf{monoton Steigend}, wenn $a_{n+1} \ge a_n \forall n \in \mathbb{N}$
        \item \textbf{streng Steigend}  $a_{n+1} > a_n \forall n \in \mathbb{N}$
            \subitem $a_n = 3n$
        \item \textbf{monoton fallend}
        \item \textbf{streng fallend}
        \item \textbf{beschränkt}  wenn $\left\{a_n \mid n \in \mathbb{N} \right\}$ (\textit{Wertebereich}) oben und unten beschränkt ist.
            \subitem Bsp. $a_n = (-1)^n$  
    \end{itemize}
\end{definition}

\subsection{Grenzwerte}
Es kann für eine Folge immer ein Glied n gefunden werden, ab dem alle weitere weniger entfernt sind von einem vermuteten Grenzwert als $\epsilon$. Wenn das gilt dann ist der Grenzwert wahr. Mit der defintion kann nur eine gegebene Grenzwertvermutung geprüft werden, sie gibt keine grenzwerte an für eine Reihe.
\begin{definition}
    [Grenzwert oder Konvergenz]

\end{definition}

\begin{definition}
    [\textit{unengentliche} Grenzwert] umgekehrt zum eigentlichen Grenzwert. Gibt es zu jedem $a \in \mathbb{R} \text{ein} m \in \mathbb{R} \text{, sodass für alle} n > m$ \footnote{Für jeden Folgeglied, gibt es einen wert m}
\end{definition}
\subsection{Harmonische Folge und Reihe}
Für diese Folge Art sollte man Grenzwert kennen.

Bei einer folge 1/n wird
bei lim inf+ die korrespondierende Reihe unendlich gross, weil 1 + 1/2 (halb) + 1/3 + 1/4 + 1/6 (halb)
Beide fälle folge, die nach 0 kovergieren aber ihre Reihen unterschiedlich sind, ein mal unednlich, ein mal Konvergenz.

\begin{theorem}
    Jede konvergente Folge ist beschränkt.
\end{theorem}